\reading{CR-20}{CVA}{STRIKE FIRST}{8}

\begin{crkeybox}[Learning objectives for this reading]
\begin{enumerate}[label=\textbf{\alph*.},leftmargin=2em]
\item Explain the motivation for and the challenges of pricing counterparty risk.
\item Describe credit value adjustment (CVA).
\item Calculate CVA and CVA as a spread with no wrong-way risk, netting, or collateralization.
\item Evaluate the impact of changes in the credit spread and recovery rate assumptions on CVA.
\item Describe debt value adjustment (DVA) and bilateral CVA (BCVA).
\item Explain the differences between unilateral CVA (UCVA) and BCVA, and between unilateral DVA (UDVA) and BCVA.
\item Calculate DVA, BCVA, and BCVA as a spread.
\item Explain how netting can be incorporated into the CVA calculation.
\item Define and calculate incremental CVA and marginal CVA and explain how to convert CVA into a running spread.
\item Explain the impact of incorporating collateralization into the CVA calculation, including the impact of margin period of risk, thresholds, and initial margins.
\item Describe wrong-way risk and contrast it with right-way risk.
\item Identify examples of wrong-way risk and examples of right-way risk.
\item Discuss the impact of collateral on wrong-way risk.
\item Identify examples of wrong-way collateral.
\item Discuss the impact of wrong-way risk on central counterparties (CCPs).
\item Describe the various wrong-way modeling methods including hazard rate approaches, structural approaches, parametric approaches, and jump approaches.
\item Explain the implications of central clearing on wrong-way risk.
\end{enumerate}
\end{crkeybox}

\begin{crtrapbox}[Seventeen objectives, not sixteen]
CR-20 is the largest reading in Credit Risk. The official 2026 document lists
\textbf{seventeen} objectives, a to q. Objective \textbf{g}, ``Calculate DVA,
BCVA, and BCVA as a spread'', is the one most often dropped, which shifts every
letter after it by one. The source notes independently letter this reading 37.a
through 37.q, confirming the count.

\smallskip
CR-20 also exists only in the Crash layer.
\end{crtrapbox}

% =====================================================================
\lo{CR-20 a}{Explain the motivation for and the challenges of pricing counterparty risk.}

\begin{itemize}
\item Accurate pricing allows financial institutions to set aside \term{reserves}
to cover potential losses if a counterparty defaults on its obligations.
\item Pricing counterparty risk helps assess the overall \term{risk profile} of a
financial institution's portfolio, enabling better risk management strategies.
\end{itemize}

\subsection*{The challenges}

\begin{enumerate}
\item \term{Exotic products.} These often require complex valuation techniques
like Monte Carlo simulation, so simplifications may be necessary for practical CVA
calculations.
\item \term{Path dependency.} Accurate future exposure assessment requires
knowledge of the \emph{entire price path} of the underlying asset, so
approximations in probability calculations are often used.
\end{enumerate}

% =====================================================================
\lo{CR-20 b}{Describe credit value adjustment (CVA).}

\begin{crdefbox}[CVA]
When taking counterparty risk into account, credit value adjustment is the
difference between the \term{risk-free value} and the \term{true value} of a
transaction or portfolio.
\end{crdefbox}

% =====================================================================
\lo{CR-20 c}{Calculate CVA and CVA as a spread with no wrong-way risk, netting, or collateralization.}

Assuming independence between PD, exposure and recovery --- that is, ignoring
wrong-way risk --- the simplified CVA expression is as follows.

\begin{crfmlbox}[Unilateral CVA]
\[ \mathrm{UCVA} \;=\; -\mathrm{LGD}\sum_{i=1}^{m}
   \mathrm{EPE}(t_i) \times \mathrm{PD}(t_{i-1},\,t_i) \]
\vk{$\mathrm{EPE}(t_i)$ is the \emph{discounted} EPE for the relevant dates in the
future; $\mathrm{PD}(t_{i-1}, t_i)$ is the marginal default probability over each
interval. The leading minus sign makes CVA a \emph{negative} number: a reduction
in the value of the portfolio.}
\end{crfmlbox}

\subsection*{CVA as a spread}

Instead of computing the CVA as a stand-alone value, it can be expressed as a
running spread, that is, a per annum charge. The formula assumes the EPE is
\emph{constant over time} and equal to the average EPE.

\begin{crfmlbox}[Unilateral CVA as a running spread]
\[ \frac{\mathrm{CVA}(t,T)}{\mathrm{CDS}_{\text{premium}}(t,T)}
   \;=\; -X^{\mathrm{CDS}} \times \text{average EPE} \]
\vk{$X^{\mathrm{CDS}}$ is the counterparty's CDS spread. The whole
LGD-and-default-probability apparatus collapses into the credit spread, which is
why this form is so much quicker to apply --- and why it is only valid under the
constant-EPE assumption.}
\end{crfmlbox}

% =====================================================================
\lo{CR-20 d}{Evaluate the impact of changes in the credit spread and recovery rate assumptions on CVA.}

\subsection*{Impact of the credit spread}

The \term{shape of the credit curve} can have a material impact on CVA.

\begin{center}
\begin{tikzpicture}
\begin{axis}[
  width=7.6cm, height=4.8cm, name=lc,
  xlabel={Maturity (years)}, ylabel={Credit spread (bps)},
  xlabel style={font=\sffamily\scriptsize}, ylabel style={font=\sffamily\scriptsize},
  tick label style={font=\sffamily\scriptsize},
  xmin=0, xmax=10, ymin=0, ymax=280, xtick={0,5,10},
  grid=major, grid style={FRMRule!50, line width=0.3pt},
  legend style={font=\scriptsize\sffamily, draw=FRMRule, fill=white,
    at={(0.5,-0.34)}, anchor=north, legend columns=-1, column sep=6pt}]
\addplot[draw=FRMRed, line width=1.1pt, smooth, domain=0.5:10, samples=40]
  {60+18*x};
\addlegendentry{upwards}
\addplot[draw=FRMNavy, line width=1.1pt, sharp plot]
  coordinates {(0.5,150) (10,150)};
\addlegendentry{flat}
\addplot[draw=FRMGreen, line width=1.1pt, smooth, domain=0.5:10, samples=40]
  {260-22*x};
\addlegendentry{inverted}
\draw[FRMGrey, dashed, line width=0.7pt] (axis cs:5,0) -- (axis cs:5,150);
\node[font=\sffamily\scriptsize, text=FRMGrey, fill=white, inner sep=1.5pt,
  align=center] at (axis cs:5.0,208) {all pinned at\\150 bps at 5y};
\end{axis}
\begin{axis}[
  ybar, bar width=15pt, width=6.8cm, height=4.8cm,
  at={(lc.right of south east)}, anchor=left of south west, xshift=1.3cm,
  ylabel={CVA}, ylabel style={font=\sffamily\scriptsize},
  tick label style={font=\sffamily\scriptsize},
  symbolic x coords={upwards,flat,inverted}, xtick=data,
  ymin=-28, ymax=0, axis x line*=top, axis y line*=left,
  nodes near coords, every node near coord/.append style={
    font=\sffamily\scriptsize\bfseries, anchor=north,
    /pgf/number format/.cd, fixed, zerofill, precision=1}]
\addplot[draw=FRMNavy, fill=FRMNavy!45, area legend]
  coordinates {(upwards,-24.6) (flat,-20.0) (inverted,-15.7)};
\end{axis}
\end{tikzpicture}
\end{center}
\figcap{Credit curve shapes on the left (illustrative, all pinned at the notes'
five-year spread of 150 bps), and the resulting CVA on a ten-year receive-fixed
interest rate swap on the right, at an LGD of 60\%. The upwards-sloping curve
gives the largest CVA, mainly because it produces the largest \emph{extrapolated
ten-year} spread; the inverted curve gives the smallest for the opposite reason.
The spread is 23\% above the flat case and 21\% below it, from curves that agree
exactly at five years.}

\begin{itemize}
\item CVA generally \term{increases} with an increasing credit spread.
\item In default, CVA will \term{converge to the exposure} of the transaction or
portfolio, which may be zero.
\end{itemize}

\subsection*{Impact of the recovery rate}

\begin{crtrapbox}[The recovery rate pulls in two directions at once]
Increasing the recovery rate will \textbf{increase} the implied probability of
default but \textbf{reduce} the resulting CVA.

\smallskip
Why both: the implied hazard rate is $\lambda = s/(1-\mathrm{RR})$, so a higher RR
shrinks the denominator and raises the implied PD. But CVA is scaled by
$\mathrm{LGD} = 1-\mathrm{RR}$, which falls faster. At a 2\% spread, moving RR
from 40\% to 60\% raises the implied hazard rate from 3.33\% to 5.00\% while
cutting LGD from 60\% to 40\%, and the LGD effect dominates.
\end{crtrapbox}

% =====================================================================
\lo{CR-20 e}{Describe debt value adjustment (DVA) and bilateral CVA (BCVA).}

\begin{crdefbox}[DVA and BCVA]
The key assumption behind unilateral CVA is that the party making the calculation
\emph{itself} cannot default. \term{Debt value adjustment (DVA)} represents
counterparty risk from the point of view of a party's \emph{own} default.

\smallskip
Consideration of a party's own default together with that of its counterparty
leads to \term{bilateral CVA (BCVA)}.
\end{crdefbox}

\begin{crfmlbox}[BCVA and its two components]
\[ \mathrm{BCVA} \;=\; \mathrm{CVA} + \mathrm{DVA} \]
\[ \mathrm{CVA} = -\mathrm{LGD}_C \sum_{i=1}^{m}
   \mathrm{EPE}(t_i)\times \mathrm{PD}_C(t_{i-1},t_i) \]
\[ \mathrm{DVA} = -\mathrm{LGD}_I \sum_{i=1}^{m}
   \mathrm{ENE}(t_i)\times \mathrm{PD}_I(t_{i-1},t_i) \]
\vk{Subscript $C$ is the counterparty and $I$ the institution itself.
$\mathrm{ENE}$ is the expected \emph{negative} exposure, that is, the EPE from the
counterparty's perspective.}
\end{crfmlbox}

\begin{crtrapbox}[Why DVA is positive when both formulas start with a minus]
Both CVA and DVA carry a leading minus sign, yet CVA is a cost and DVA is a
benefit. The resolution is the sign of ENE.

\smallskip
ENE is \emph{negative} by construction, so $-\mathrm{LGD}_I \times
(\text{negative}) \times \mathrm{PD}_I$ comes out \textbf{positive}. Worked
through with $\mathrm{EPE}=100$, $\mathrm{ENE}=-80$, both LGDs 60\%,
$\mathrm{PD}_C = 2\%$ and $\mathrm{PD}_I = 1.5\%$:
\[ \mathrm{CVA} = -0.6(100)(0.02) = -1.20 \qquad
   \mathrm{DVA} = -0.6(-80)(0.015) = +0.72 \]
\[ \mathrm{BCVA} = -1.20 + 0.72 = -0.48 \]
This is the same sign convention used at \textbf{CR-12 b}, where the
credit-adjusted value is $f_{nd} - \mathrm{CVA} + \mathrm{DVA}$, and at
\textbf{CR-15 h}, where the xVA table records CVA as negative and DVA as positive.
\end{crtrapbox}

% =====================================================================
\lo{CR-20 f}{Explain the differences between unilateral CVA (UCVA) and BCVA, and between unilateral DVA (UDVA) and BCVA.}

\begin{center}
\footnotesize
\begin{tabular}{L{2.6cm} L{5.7cm} L{6.2cm}}
\toprule
\thead{Measure} & \thead{Whose default is considered} & \thead{Relation to BCVA} \\
\midrule
\term{UCVA} & The \emph{counterparty's} default only. The institution assumes it
  cannot itself default &
  BCVA is UCVA \emph{plus} the DVA term, so BCVA is less negative than UCVA \\
\term{UDVA} & The \emph{institution's own} default only &
  BCVA is UDVA \emph{plus} the CVA term \\
\term{BCVA} & \emph{Both} parties' defaults, each weighted by the other's
  survival &
  --- \\
\bottomrule
\end{tabular}
\end{center}
\figcap{The distinction is entirely about whose default is admitted into the
calculation. A unilateral measure switches one of the two terms off.}

% =====================================================================
\lo{CR-20 g}{Calculate DVA, BCVA, and BCVA as a spread.}

The DVA and BCVA calculations are the formulas at objective e, worked through in
the trap box there. The spread form extends the CVA spread of objective c to both
sides.

\begin{crfmlbox}[BCVA as a running spread]
\[ \frac{\mathrm{BCVA}(t,T)}{\mathrm{CDS}_{\text{premium}}(t,T)}
   \;=\; -X^{\mathrm{CDS}}_{C}\times \text{average EPE}
   \;-\; X^{\mathrm{CDS}}_{I}\times \text{average ENE} \]
\vk{$X^{\mathrm{CDS}}_{C}$ and $X^{\mathrm{CDS}}_{I}$ are the CDS spreads of the
counterparty and of the institution itself. The second term is positive in value
because average ENE is negative, exactly as in the stand-alone case.}
\end{crfmlbox}

\begin{crkeybox}[The structure to hold]
Each party's \emph{own credit spread} multiplies the \emph{other side's}
exposure. The counterparty's spread is applied to your positive exposure; your own
spread is applied to your negative exposure. Swapping which spread goes with which
exposure is the single easiest error in this formula.
\end{crkeybox}

% =====================================================================
\lo{CR-20 h}{Explain how netting can be incorporated into the CVA calculation.}

Netting reduces the CVA price, as it nets --- that is, reduces --- exposure when
trades are settled. Because CVA is linear in EPE, anything that lowers EPE lowers
CVA in the same proportion.

% =====================================================================
\lo{CR-20 i}{Define and calculate incremental CVA and marginal CVA and explain how to convert CVA into a running spread.}

\subsection*{Incremental CVA}

Incremental CVA measures the change in CVA caused by a new trade, considering the
impact of netting with existing trades. It helps assess the impact of a new trade
on the overall CVA risk profile.

\begin{crfmlbox}[Incremental CVA]
\[ \Delta\mathrm{CVA}_{NS \to NS'} \;=\; \mathrm{CVA}_{NS'} - \mathrm{CVA}_{NS} \]
\[ \Delta\mathrm{CVA}_{NS \to NS'} \;=\; -\mathrm{LGD}\sum_{i=1}^{m}
   \mathrm{EPE}_{NS \to NS'}(t_i)\times \mathrm{PD}(t_{i-1},t_i) \]
\vk{$\mathrm{CVA}_{NS}$ and $\mathrm{CVA}_{NS'}$ are the portfolio CVA before and
after the change. The second line is the standalone formula of objective c with
the \emph{incremental} EPE replacing the standalone EPE. This assumes the change
does not alter the PD or LGD of the counterparty.}
\end{crfmlbox}

\begin{crtrapbox}[Correction: the source notes set two different things equal]
The notes print this as ``$V(i) = \Delta\mathrm{CVA}_{NS,i} = \mathrm{CVA}(NS,i)
- \mathrm{CVA}(NS)$'' and then define, in the same variable key, $V(i)$ as ``the
risk-free value of new trade $i$''. Both cannot hold: the risk-free value of a new
trade is not its incremental CVA.

\smallskip
The GARP form carries no $V(i)$ term. The incremental CVA is the portfolio CVA
after the change less the portfolio CVA before it, and nothing else.
\end{crtrapbox}

Incremental CVA ensures that the CVA of a new transaction is represented by its
contribution to the overall CVA at the time it is executed, which is what makes it
chargeable to individual traders, businesses and ultimately clients. It applies
not only to new transactions but also to unwinding, restructuring, and changes to
contractual terms such as margin agreements.

\begin{crkeybox}[Order matters, but only once]
CVA depends on the \emph{order} in which transactions are executed, but does not
change due to subsequent transactions. A trade charged at execution keeps its
charge.
\end{crkeybox}

\subsection*{Marginal CVA}

Marginal CVA isolates the CVA contribution of individual trades \emph{within} a
netted portfolio. It provides a detailed analysis of which trades have the biggest
impact on a counterparty's CVA. The formula is the same as standalone CVA but uses
the \term{marginal expected exposure} instead of the standalone exposure.

\begin{center}
\footnotesize
\begin{tabular}{L{2.8cm} L{5.7cm} L{6.0cm}}
\toprule
 & \thead{Incremental CVA} & \thead{Marginal CVA} \\
\midrule
Question answered & What does \emph{adding this trade} do to portfolio CVA? &
  How is the \emph{existing} portfolio CVA attributed across its trades? \\
Exposure used & Incremental EPE, i.e.\ the change in EPE &
  Marginal EPE \\
When used & At execution, to charge the new trade &
  After the fact, to analyse an existing book \\
\bottomrule
\end{tabular}
\end{center}
\figcap{Both formulas are the standalone CVA with a different exposure measure
substituted in. What separates them is the question, not the mathematics:
incremental is forward-looking and trade-by-trade at execution, marginal is a
decomposition of a book that already exists.}

\subsection*{Converting CVA into a running spread}

To integrate CVA costs into the pricing of instruments like swaps, divide the
standalone CVA by the product of the risky duration and the notional amount.

\begin{crexbox}[CVA as a running spread on a swap]
For a CVA of $-90{,}000$ and a swap with a notional of \$100 million and a risky
duration of 3.75, what is the running spread?
\tcblower
\[ \frac{-90{,}000}{3.75 \times \$100{,}000{,}000}
   = \frac{-90{,}000}{375{,}000{,}000}
   = -0.000240 \]
which is \textbf{$-2.40$ basis points}.

\smallskip
Check the direction: the CVA is negative and so is the spread, meaning the charge
is levied \emph{against} the counterparty in the pricing of the swap.
\end{crexbox}

% =====================================================================
\lo{CR-20 j}{Explain the impact of incorporating collateralization into the CVA calculation, including the impact of margin period of risk, thresholds, and initial margins.}

\begin{center}
\footnotesize
\begin{tabular}{L{4.4cm} L{10.1cm}}
\toprule
\thead{Feature} & \thead{Effect on CVA} \\
\midrule
Collateral generally & \textbf{Reduces} CVA by lowering the counterparty's
  expected exposure. It does \emph{not} affect the default probability \\
Minimum transfer amounts & \textbf{Increase} CVA \\
Thresholds & \textbf{Increase} CVA \\
Higher initial margins \newline (negative thresholds) & \textbf{Decrease} CVA \\
Margin period of risk & Defines the timeframe over which CVA is measured. As the
  MPoR increases, CVA gradually approaches the \emph{uncollateralized} CVA \\
\bottomrule
\end{tabular}
\end{center}
\figcap{The directions match CR-19 h exactly: lower MPoR, threshold and MTA are
better, higher initial margin is better. Note the mechanism stated in the first
row --- collateral works on \emph{exposure}, never on the probability of default.}

% =====================================================================
\lo{CR-20 k}{Describe wrong-way risk and contrast it with right-way risk.}

\begin{center}
\footnotesize
\begin{tabular}{L{3.0cm} L{5.5cm} L{6.0cm}}
\toprule
 & \thead{Wrong-way risk (WWR)} & \thead{Right-way risk (RWR)} \\
\midrule
What happens & Exposure \textbf{increases} as the counterparty's credit quality
  \textbf{worsens} &
  Exposure \textbf{decreases} as the counterparty's credit quality
  \textbf{improves} \\
Effect on CVA & \textbf{Increases} CVA & \textbf{Reduces} CVA \\
Effect on DVA & \textbf{Reduces} DVA & \textbf{Increases} DVA \\
\bottomrule
\end{tabular}
\end{center}
\figcap{Three directions per column and all three flip together. The underlying
idea is a single correlation: wrong-way risk is positive correlation between your
exposure and the counterparty's default probability.}

% =====================================================================
\lo{CR-20 l}{Identify examples of wrong-way risk and examples of right-way risk.}

\begin{center}
\footnotesize
\begin{tabular}{L{3.3cm} L{11.2cm}}
\toprule
\thead{Transaction} & \thead{Why it is wrong-way} \\
\midrule
\term{OTC put option} & You buy an out-of-the-money put on a stock from the
  issuer. If the stock price falls (good for you), the issuer's creditworthiness
  might also decline (bad for you). \textbf{Out-of-the-money puts have more WWR
  than in-the-money puts} \\
\term{Credit default swap} & You buy CDS protection against defaults on bonds
  issued by a reference entity. If the reference entity defaults the value of
  your protection increases, but the gain on the CDS and the default probability
  of the \emph{protection seller} rise together \\
\term{Foreign currency transactions} & You lend currency to a counterparty in an
  emerging market. If their currency weakens (good for you), their
  creditworthiness might also deteriorate \\
\term{Foreign currency swaps} & You swap fixed-rate payments with a counterparty
  in another currency. If interest rates fall in their currency (good for you),
  their creditworthiness might also decline \\
\term{Interest rate swaps} & You receive fixed and pay floating. If rates rise
  (good for you), their creditworthiness might also decline \\
\term{Commodity forwards} & You buy a forward to lock in a low oil price. If oil
  prices rise (good for you), the counterparty might face pressure due to
  concentrated positions \\
\bottomrule
\end{tabular}
\end{center}
\figcap{Every row has the same shape: the move that is good for you is the move
that is bad for them. That is the test to apply to any candidate example, and it
is also why right-way risk examples are simply the same trades with the
correlation reversed.}

% =====================================================================
\lo{CR-20 m}{Discuss the impact of collateral on wrong-way risk.}

Collateral reduces exposure, but its effectiveness depends on \term{timing}.

\begin{itemize}
\item \term{Gradual exposure increase.} Collateral helps mitigate WWR.
\item \term{Sudden exposure jump}, for example a currency devaluation. Collateral
might be difficult to obtain in time.
\end{itemize}

\begin{crtrapbox}[Collateral is weakest exactly where wrong-way risk is worst]
The wrong-way exposures that matter most are the jump-type ones --- a
devaluation, a default of a reference entity --- and those are precisely the ones
where the margin period of risk means the collateral call cannot be met in time.
\end{crtrapbox}

% =====================================================================
\lo{CR-20 n}{Identify examples of wrong-way collateral.}

\begin{crkeybox}[Wrong-way collateral]
A \term{payer interest rate swap collateralized by government bonds}. When
interest rates rise you benefit on the swap, so your exposure grows --- but the
value of the government bonds posted as collateral falls at the same moment.

\smallskip
The exposure and the collateral move against each other, which is wrong-way risk
located \emph{inside the collateral} rather than in the counterparty.
\end{crkeybox}

% =====================================================================
\lo{CR-20 o}{Discuss the impact of wrong-way risk on central counterparties (CCPs).}

CCPs rely on collateral and default funds from members to cover losses. WWR makes
CCPs vulnerable in two ways.

\begin{enumerate}[label=\alph*)]
\item \term{Higher credit quality members may have higher WWR}, requiring
adjustments in margin and default fund contributions.
\item \term{Risky collateral posted by members} can increase WWR for the CCP.
\end{enumerate}

CCPs are therefore susceptible to WWR and need adjustments based on member credit
quality and collateral risk.

% =====================================================================
\lo{CR-20 p}{Describe the various wrong-way modeling methods including hazard rate approaches, structural approaches, parametric approaches, and jump approaches.}

\begin{center}
\footnotesize
\begin{tabular}{L{3.0cm} L{5.6cm} L{5.9cm}}
\toprule
\thead{Method} & \thead{How it works} & \thead{Assessment} \\
\midrule
\term{Hazard rate} \newline (intensity approach) &
  Simulates credit spreads and calculates expected exposure when default occurs,
  that is, at wider spreads &
  Easy to implement, but \textbf{may underestimate WWR} due to weak dependence
  between default and exposure \\
\term{Structural} &
  Maps default and exposure distributions from a \emph{bivariate} distribution.
  WWR is identified as early default with high exposure, that is, positive
  correlation &
  \emph{Advantage}: uses the existing exposure distribution.
  \emph{Disadvantage}: the information might not be relevant \\
\term{Parametric} &
  Examines the historical link between portfolio exposure and credit spreads. High
  portfolio value together with a high credit spread indicates higher WWR &
  Relies on historical data reflecting current scenarios \\
\term{Jump} &
  Assumes a \emph{jump} at default, for example a sharp currency devaluation. The
  jump factor reflects the size of the devaluation &
  \textbf{Most applicable to WWR} \\
\bottomrule
\end{tabular}
\end{center}
\figcap{The two ends of the table are the examinable pair. The hazard rate
approach is the easiest to implement and the one most likely to \emph{under}state
wrong-way risk; the jump approach is the one best suited to it, because the
wrong-way events that matter are discontinuous.}

% =====================================================================
\lo{CR-20 q}{Explain the implications of central clearing on wrong-way risk.}

Central clearing does not remove wrong-way risk; it relocates it. Because a CCP's
protection is built on collateral and default fund contributions, and both of
those can be impaired by the same events that raise member exposures, the CCP
inherits the wrong-way risk of its members' positions and of the collateral they
post.

The two adjustments a CCP must therefore make are the ones at objective o: margin
and default fund contributions scaled to member credit quality, and haircuts or
eligibility rules scaled to the wrong-way risk in the collateral itself.
